\subsubsection{The new Form of \bsq{switch}}\label{sec:NewSwitch}
The alternative syntax for \lstinline|switch|\index{switch} comes in two flavours (as \textit{statement}\index{switch!statement} and \textit{expression}\index{switch!expression}).

The new \lstinline|switch|\index{switch!statement} statement looks like this:

\begin{lstlisting}[numbers=left]
// switch statement
switch( <selector> )
{
    case <switchlabel1> -> <singleStatement>;

    case <switchlabel2>, <switchlabel3> -> <singleStatement>;

    case <switchlabel4> ->
    {
        <statements>;
    }

    case <switchlabel4>, <switchlabel5> ->
    {
        <statements>;
    }

    default -> <singleStatement>;
}
\end{lstlisting}

And this is the new \lstinline|switch|\index{switch!expression} expression:

\begin{lstlisting}[numbers=left]
// switch expression
var result = switch( <selector> )
{
    case <switchlabel1> -> <expression>;

    case <switchlabel2>, <switchlabel3> -> <expression>;

    case <switchlabel4> ->
    {
        <statements>;
        yield <value>;
    }

    case <switchlabel4>, <switchlabel5> ->
    {
        <statements>;
        yield <value>;
    }

    default -> <expression>;
}
\end{lstlisting}

Of course the \lstinline|default| in line~18 can be followed by a statement block as the \lstinline|case| in line~8, and the \lstinline|default| in line~40 can have a complex expression like the \lstinline|case| in line~28.

As this format does not allow a fall-through in cases where the branches for two or more labels are doing the same stuff, it is possible here to have more than one switch label per \lstinline|case|, as shown in lines~6, 13, 26 and 34.

If used with pattern matching (refer to \autocite{ORACLE_DOC_PATTERNMATCHING}), a \lstinline|switch| looks like this:
\begin{lstlisting}[numbers=left]
var selector = <anObject>
// switch statement
switch( selector )
{
    case null -> <singleStatement>;
    
    case <class1> <name1> -> <singleStatement>;

    case <class2> <name2> ->
    {
        <statements>;
    }

    default -> <singleStatement>;
}

// switch statement
var result = switch( selector )
{
    case null -> <expression>;

    case <class1> <name1> -> <expression>;

    case <class2> <name2> ->
    {
        <expressions>;
        yield <value>;
    }

    default -> <expression>;
}
\end{lstlisting}

Same as for the \lstinline|default| branch, the \lstinline|case null| branch can be a statement block or a complex expression.

And an expression can also be always a \lstinline|throw|.


%==============================================================================
% End of file!!
%==============================================================================
